% !TEX program = uplatex

\documentclass[uplatex,dvipdfmx,a4paper,11pt]{jsarticle}

\usepackage[top=25mm,bottom=25mm,left=22mm,right=22mm]{geometry}

\usepackage{longtable,booktabs,array,tabularx}

\usepackage[dvipdfmx]{xcolor,hyperref}

\usepackage{pxjahyper}

\usepackage{enumitem}

\hypersetup{colorlinks=true,linkcolor=black,urlcolor=blue}

\setlength{\parindent}{0pt}

\setlength{\parskip}{0.65em}

\renewcommand{\arraystretch}{1.25}

\title{14章 ソフトウェア工学専門職実務\\\large Software Engineering Professional Practice の日本語まとめ}

\author{SWEBOK Guide v4.0a Chapter 14 を初学者向けに再構成}

\date{作成日：2026年5月6日}

\begin{document}

\maketitle

\begin{center}\small 本資料はSWEBOK Guide v4.0a Chapter 14を初学者向けに要約・再構成したものであり、逐語訳ではない。\end{center}

\section*{この資料の主張}

ソフトウェア工学専門職実務は、技術だけでなく、倫理、標準、法的責任、契約、文書化、チームでの協働、利害関係者との対話、読み書きと発表を含む実務能力である。ソフトウェアが社会の中で広く使われるほど、技術者の判断には説明責任と公共への配慮が求められる。

\section*{この資料の読み方}

専門用語は原則として日本語で示し、必要に応じて英語を括弧内に併記する。英語名は、元資料やツール上の名称を確認したいときの手がかりとして残す。初学者が前提知識なしに読めるように、各節では「何のための概念か」「実務では何を見ればよいか」を重視する。

\tableofcontents

\newpage

\section*{全体概要：この章で学ぶこと}

\addcontentsline{toc}{section}{全体概要：この章で学ぶこと}

この章は、ソフトウェア工学を専門職として実践するために必要な知識、技能、態度を扱う。ここでいう専門職実務とは、単にプログラムを書く能力ではない。倫理、標準、法的論点、契約、文書化、チームでの働き方、利害関係者との関係、読み書きや発表の技能を含む総合的な実務能力である。

ソフトウェアは、金融、医療、交通、教育、行政、娯楽、家庭内機器など、社会生活の広い範囲に入り込んでいる。したがって、ソフトウェア技術者の判断は、個人の利便性だけでなく、安全、プライバシー、財産、組織活動、社会的信頼にも影響する。第14章は、このような影響を理解し、責任ある態度で仕事を進めるための土台を示す章である。

\begin{center}
\fbox{\parbox{0.94\linewidth}{\textbf{到達目標}\par この章を読み終えると、専門職性、倫理規程、専門職団体、標準、契約、法的論点、文書化、トレードオフ分析、チーム力学、利害関係者との関係、コミュニケーション技能の役割を説明できるようになる。}}
\end{center}

\begin{center}
\fbox{\parbox{0.94\linewidth}{\textbf{TODO（図） 14章の全体地図を入れる。}\par\textbf{画像生成用プロンプト}\par 白背景、モノクロ、A4本文幅に収まる横長の学習用図解を作成する。中央に「ソフトウェア工学専門職実務」を置き、三つの大分類として「専門性」「集団力学と心理学」「コミュニケーション技能」を配置する。「専門性」から「倫理」「標準」「契約」「法的論点」「文書化」「トレードオフ分析」へ枝分かれさせる。日本語ラベル、細線、講義ノート風、余白を広めにする。}}
\end{center}

\section*{最初に必要な用語}

\addcontentsline{toc}{section}{最初に必要な用語}

\begin{center}
\small
\begin{longtable}{p{0.24\linewidth}p{0.71\linewidth}}
\toprule
\textbf{用語} & \textbf{初学者向けの説明} \\
\midrule
\endfirsthead
\toprule
\textbf{用語} & \textbf{初学者向けの説明} \\
\midrule
\endhead
専門職実務 & 社会的責任を伴う仕事を、一定の標準・倫理・技能に従って行うこと。 \\
\midrule
専門性 & 専門分野の知識だけでなく、責任ある判断、継続学習、職業倫理を含む態度である。 \\
\midrule
倫理規程 & 専門家が守るべき価値観と行動原則を明文化したもの。 \\
\midrule
専門職行動規範 & 専門家として実際に避けるべき行為、行うべき行為を示す規範である。 \\
\midrule
専門職団体 & 専門分野の知識体系、標準、資格、倫理、教育、交流を支える団体である。 \\
\midrule
標準 & 成果物やプロセスが最低限満たすべき性質や手順を定めた合意文書である。 \\
\midrule
認定 & 教育機関やプログラムが一定の基準を満たすことを第三者が認めること。 \\
\midrule
資格証明 & 個人が特定分野の知識・能力を持つことを試験などで確認すること。 \\
\midrule
免許 & 特定の業務を行う権限を、政府または法定機関が個人に与えること。 \\
\midrule
知的財産 & 発明、著作物、設計、商標、営業秘密など、知的活動から生まれる財産的価値である。 \\
\midrule
利害関係者 & ソフトウェアの開発・運用結果に影響を与える、または影響を受ける人や組織である。 \\
\midrule
トレードオフ & 一方を良くすると他方が悪くなる関係の中で、目的に照らして釣り合いを取る判断である。 \\
\midrule
集団力学 & チームや組織の中で、人の行動、協力、衝突、意思決定がどのように生じるかを扱う考え方である。 \\
\midrule
包摂性 & 性別、文化、言語、地域、障害、経験などの違いを前提に、人が参加しやすい状態を作ること。 \\
\midrule
\bottomrule
\end{longtable}
\normalsize
\end{center}

\section*{略語}

\addcontentsline{toc}{section}{略語}

\begin{center}
\small
\begin{longtable}{p{0.17\linewidth}p{0.78\linewidth}}
\toprule
\textbf{略語} & \textbf{日本語での意味} \\
\midrule
\endfirsthead
\toprule
\textbf{略語} & \textbf{日本語での意味} \\
\midrule
\endhead
ACM & 計算機協会。情報処理分野の国際的な学会・専門職団体である。 \\
\midrule
CCPA & カリフォルニア州消費者プライバシー法。データプライバシー法の例である。 \\
\midrule
EEA & 欧州経済領域。 \\
\midrule
ENAEE & 欧州工学教育認定ネットワーク。工学教育認定に関わる組織である。 \\
\midrule
EU & 欧州連合。 \\
\midrule
GDPR & 一般データ保護規則。EU・EEAにおける個人データ保護規則の代表例である。 \\
\midrule
IEA & 国際エンジニアリング連合。工学教育・専門職認定に関わる国際的枠組みである。 \\
\midrule
IEC & 国際電気標準会議。電気・電子・関連技術の国際標準化組織である。 \\
\midrule
IEEE CS & IEEEコンピュータソサエティ。 \\
\midrule
IFIP & 国際情報処理連合。 \\
\midrule
IP & 知的財産。 \\
\midrule
ISO & 国際標準化機構。 \\
\midrule
NDA & 秘密保持契約。仕事を通じて知った機密情報を守る契約である。 \\
\midrule
UI/UX & 利用者インタフェース／利用者体験。画面や操作感だけでなく、利用者が受け取る体験を含む。 \\
\midrule
WIPO & 世界知的所有権機関。 \\
\midrule
WTO & 世界貿易機関。 \\
\midrule
\bottomrule
\end{longtable}
\normalsize
\end{center}

\section*{導入：専門職実務が重要な理由}

\addcontentsline{toc}{section}{導入：専門職実務が重要な理由}

ソフトウェア工学専門職実務は、ソフトウェア技術者が専門的、責任ある、倫理的な方法で仕事をするために必要な知識、技能、態度を扱う知識領域である。専門職実務は、成果物だけでなく、作業の過程にも基準を求める。つまり「動くものを作ったからよい」ではなく、「どのような根拠、手順、責任、説明可能性をもって作ったか」も問われる。

専門職としての実践には、一般に受け入れられた知識体系、倫理規程と行動規範、認定・資格証明・免許の仕組み、専門職団体、標準が関わる。ソフトウェア技術者は、ACM、IEEE、IFIP、IEEEコンピュータソサエティ、ISO/IEC、ISO/IEC/IEEEなどが示す標準や指針を参考にしながら、自分の仕事を説明できる形で進める必要がある。

第14章は三つの大領域で構成される。第一は専門性であり、倫理、標準、法的論点、契約、文書化、トレードオフ分析を扱う。第二は集団力学と心理学であり、チーム作業、個人の認知、複雑な問題、利害関係者、不確実性、多様性を扱う。第三はコミュニケーション技能であり、読む、書く、チーム内で伝える、発表する能力を扱う。

\section*{1 専門性}

\addcontentsline{toc}{section}{1 専門性}

専門性とは、専門分野の知識を持っているだけでなく、その知識を社会的責任のもとで使う態度を含む。ソフトウェア技術者の専門性は、倫理規程と専門職行動規範に従うこと、標準と実務慣行を理解すること、継続的に学ぶこと、利害関係者へ説明できる判断を行うことに表れる。

専門職団体は、知識体系、倫理規程、専門職行動規範、認定や免許の基準を示す。これらの基準は、技術者が共同体の一員として認められるための入口であると同時に、能力や過失を判断する基準にもなり得る。

\subsection*{1.1 認定・資格証明・資格・免許}

\addcontentsline{toc}{subsection}{1.1 認定・資格証明・資格・免許}

認定、資格証明、資格、免許は似ているが、対象と強制力が異なる。初学者は「組織を見るのか、個人を見るのか」「任意なのか、法的に必要なのか」で区別するとよい。

\begin{center}
\small
\begin{longtable}{p{0.15\linewidth}p{0.19\linewidth}p{0.33\linewidth}p{0.28\linewidth}}
\toprule
\textbf{区分} & \textbf{対象} & \textbf{意味} & \textbf{実務上の見方} \\
\midrule
\endfirsthead
\toprule
\textbf{区分} & \textbf{対象} & \textbf{意味} & \textbf{実務上の見方} \\
\midrule
\endhead
認定 & 教育機関・教育課程・組織 & 一定の教育・品質基準を満たしていることを第三者が認める。 & 工学教育プログラムや大学課程の品質を示す。 \\
\midrule
資格証明 & 個人 & 試験や経験要件により、一定水準の知識・能力を持つことを示す。 & 能力を外部に説明する材料になるが、多くの場合は任意である。 \\
\midrule
資格 & 個人 & 特定の能力や条件を満たしている状態を示す。再資格化が不要な場合もある。 & 資格証明に近いが、制度により意味が異なる。 \\
\midrule
免許 & 個人 & 特定の業務を行う権限を政府や法定機関が与える。 & 国や地域によっては、専門技術者として業務を行う条件になる。 \\
\midrule
\bottomrule
\end{longtable}
\normalsize
\end{center}

\begin{center}
\fbox{\parbox{0.94\linewidth}{\textbf{注意}\par SWEBOK第14章では、ソフトウェア工学における資格証明や資格は多くの場合任意であり、資格がないことだけでソフトウェア技術者として働けないわけではないと説明している。一方で、免許は国や地域の制度に依存し、業務権限に関わる場合がある。}}
\end{center}

\begin{center}
\fbox{\parbox{0.94\linewidth}{\textbf{TODO（図） 認定・資格証明・免許の違いを示す比較図を入れる。}\par\textbf{画像生成用プロンプト}\par 白背景、モノクロ、A4本文幅の比較図を作成する。横軸に「組織を対象」から「個人を対象」、縦軸に「任意性が高い」から「法的強制力が高い」を置く。点として「認定」「資格証明」「資格」「免許」を配置し、それぞれに短い説明を添える。日本語ラベル、講義ノート風。}}
\end{center}

\subsubsection*{1.1.1 認定}

\addcontentsline{toc}{subsubsection}{1.1.1 認定}

認定とは、組織、教育機関、教育プログラムなどが一定の能力、権限、信頼性、品質を備えていることを認めることである。工学分野では、認定された教育課程を修了することが、技術者が基礎知識を得る主要な方法になる国もある。国際的な枠組みとして、国際エンジニアリング連合や欧州工学教育認定ネットワークが例として挙げられる。

\subsubsection*{1.1.2 資格証明と資格}

\addcontentsline{toc}{subsubsection}{1.1.2 資格証明と資格}

資格証明とは、個人が特定の専門領域で一定水準の能力を持つことを確認する仕組みである。通常は試験と経験要件によって取得し、専門職団体などの非政府組織が運営することが多い。ソフトウェア工学では、標準的な実務知識を確認し、キャリア形成を支える目的がある。

資格証明は、専門職標準を満たす力、専門的判断を適用する力、所定の知識、実務経験、実証された方法論の理解を示す。ただし、それは能力を示す一つの手段であり、すべてのソフトウェア技術者が何らかの資格証明を持つわけではない。

\subsubsection*{1.1.3 免許}

\addcontentsline{toc}{subsubsection}{1.1.3 免許}

免許とは、ある人に特定の活動を行う権限を与え、その成果に責任を持たせる制度である。通常は政府機関または法定機関が発行する。免許制度の目的は、資格を持たない者による危険な実務から公衆を守ることである。国や地域によっては、免許を持たない者が専門技術者として業務を行ったり、企業が工学サービスを提供したりできない場合がある。

\subsection*{1.2 倫理規程と専門職行動規範}

\addcontentsline{toc}{subsection}{1.2 倫理規程と専門職行動規範}

倫理規程と専門職行動規範は、技術者の実務と意思決定がどのような価値と行動に基づくべきかを示す。ソフトウェア技術者は、顧客や雇用主だけでなく、公衆の健康、安全、福祉を考慮しなければならない。これは、ソフトウェアの失敗が利用者の生活、社会秩序、財産、生命に影響し得るからである。

倫理規程は、社会規範や地域の法律の文脈の中で成立する。複数の義務が衝突する場面、たとえば顧客の利益、利用者の安全、守秘義務、公共の利益がぶつかる場面で、倫理規程は判断の手がかりになる。ACM、IEEE、IFIPなどの専門職団体は、それぞれ倫理規程や専門職行動規範を公開している。

\begin{center}
\fbox{\parbox{0.94\linewidth}{\textbf{倫理違反の例}\par 倫理違反は、何かを積極的に行う作為だけでなく、必要なことをしない不作為でも起こる。たとえば、不十分な仕事を隠すこと、機密情報を漏らすこと、事実を偽ること、能力を誇張することは作為の違反である。重大なリスクを知らせないこと、参考文献や貢献者を示さないこと、顧客の正当な利益を代表しないことは不作為の違反である。}}
\end{center}

\begin{center}
\fbox{\parbox{0.94\linewidth}{\textbf{TODO（図） 倫理規程違反の種類を示す図を入れる。}\par\textbf{画像生成用プロンプト}\par 白背景、モノクロ、A4本文幅の二分割図を作成する。左に「作為による違反」、右に「不作為による違反」を配置する。左には「機密漏えい」「事実の偽造」「能力の虚偽表示」「不十分な仕事の隠蔽」を、右には「リスク未開示」「重要情報の未提供」「参考文献や貢献の未表示」「顧客利益の不十分な代表」を入れる。日本語ラベル、講義ノート風。}}
\end{center}

\subsection*{1.3 専門職団体の性質と役割}

\addcontentsline{toc}{subsection}{1.3 専門職団体の性質と役割}

専門職団体は、実務家と研究者を含む共同体である。専門職団体は、専門分野を定義し、発展させ、必要に応じて規律を与える。ソフトウェア工学では、専門職団体が知識体系、標準、倫理規程、資格、認定、研修、出版、会議を支える。

個人にとって、専門職団体への参加は、知識を更新し、専門分野とのつながりを保ち、専門家ネットワークを広げる手段である。技術の変化が速いソフトウェア分野では、専門職団体を通じた継続学習が特に重要である。

\begin{itemize}[leftmargin=2em]

\item 一般に受け入れられた知識体系を整備し、広める。

\item 免許、資格証明、認定の基盤を提供する。

\item 倫理規程や専門職行動規範を示し、違反への懲戒を行うことがある。

\item 会議、研修、出版、標準化活動を通じて専門分野を発展させる。

\end{itemize}

\subsection*{1.4 ソフトウェア工学標準の性質と役割}

\addcontentsline{toc}{subsection}{1.4 ソフトウェア工学標準の性質と役割}

ソフトウェア工学標準は、開発、保守、支援、品質、要求、設計、テスト、運用など幅広い話題を扱う。標準は、個々の組織や技術者が毎回ゼロから判断しなくてもよいように、過去の経験と合意をまとめたものである。

標準に従うことには二つの価値がある。第一に、成果物や作業の最低限の性質を明確にし、暗黙の前提や過信を減らす。第二に、法的な争いが生じたときに、一般に受け入れられた実務に従っていたことを説明する根拠になり得る。

\begin{center}
\fbox{\parbox{0.94\linewidth}{\textbf{標準は創造性の敵ではない}\par 標準は「考えなくてもよい」という意味ではない。むしろ、最低限の共通土台を置くことで、技術者が本当に考えるべき問題に集中しやすくする。}}
\end{center}

\subsection*{1.5 ソフトウェアの経済的影響}

\addcontentsline{toc}{subsection}{1.5 ソフトウェアの経済的影響}

ソフトウェアは、個人、企業、社会の各水準で経済的影響を持つ。個人の水準では、技術者の雇用や評価は、顧客や雇用主のニーズを理解し、実行できるかに左右される。顧客や雇用主の財務状況も、ソフトウェア購入や開発の成功・失敗に影響される。

企業の水準では、適切なソフトウェアは何か月分もの作業を削減し、利益や組織効率を高める。一方で、ソフトウェアの開発・取得には大きな費用がかかり、失敗すれば大きな損失となる。社会の水準では、ソフトウェアの成功や失敗は、事故、サービス停止、生産性、社会秩序、財産、生命にまで影響し得る。

\subsection*{1.6 雇用契約}

\addcontentsline{toc}{subsection}{1.6 雇用契約}

ソフトウェア工学サービスは、企業間契約、コンサルティング契約、直接雇用、ボランティアなど多様な関係のもとで提供される。契約では、何を提供するか、その対価は何か、どの条件で働くかが定められる。

契約で特に重要なのは、秘密保持と知的財産の所有である。技術者は、仕事を通じて知った情報を外部に漏らさないために、秘密保持契約や知的財産契約に署名することがある。これらの義務は、契約終了後も続く場合がある。

\begin{center}
\small
\begin{longtable}{p{0.28\linewidth}p{0.66\linewidth}}
\toprule
\textbf{契約で扱われやすい項目} & \textbf{説明} \\
\midrule
\endfirsthead
\toprule
\textbf{契約で扱われやすい項目} & \textbf{説明} \\
\midrule
\endhead
秘密保持 & 仕事を通じて知った情報、設計、顧客情報、事業情報を守る。 \\
\midrule
知的財産の所有 & 成果物、発明、発見、アイデアの権利が誰に帰属するかを決める。 \\
\midrule
作業場所 & 顧客先、自宅、開発拠点、リモートなどを定める。 \\
\midrule
適用標準 & 成果物や作業が従う標準、規程、手順を定める。 \\
\midrule
開発環境 & 使用するハードウェア、ソフトウェア、構成を定める。 \\
\midrule
責任制限 & 損害や不具合が生じた場合の責任範囲を定める。 \\
\midrule
連絡・エスカレーション & 誰に、どの順序で、どのように報告・相談するかを定める。 \\
\midrule
報酬・勤務条件 & 単価、支払頻度、勤務時間、勤務条件を定める。 \\
\midrule
\bottomrule
\end{longtable}
\normalsize
\end{center}

\subsection*{1.7 法的論点}

\addcontentsline{toc}{subsection}{1.7 法的論点}

ソフトウェア工学の専門職実務には、多くの法的論点が関わる。標準、商標、特許、著作権、営業秘密、専門職責任、法的要求、貿易遵守、サイバー犯罪、データプライバシーが代表例である。法律は国や地域に強く依存するため、実務で判断が必要な場合は、その法域と分野に詳しい弁護士などの専門家に確認する必要がある。

\begin{center}
\fbox{\parbox{0.94\linewidth}{\textbf{法的注意}\par この資料はSWEBOK第14章の学習用要約であり、法的助言ではない。実務では、契約、知的財産、輸出入、個人情報、責任範囲について、関係する国・地域の専門家に確認する必要がある。}}
\end{center}

\subsubsection*{1.7.1 標準}

\addcontentsline{toc}{subsubsection}{1.7.1 標準}

標準に従うことは、法的行動や専門過誤の主張に対する防御材料になり得る。これは、標準が一般に受け入れられた実務の根拠を示すからである。

\subsubsection*{1.7.2 商標}

\addcontentsline{toc}{subsubsection}{1.7.2 商標}

商標とは、商品やサービスの出所を示すために用いられる名称、記号、ロゴ、画像、包装などである。商標は、他者が同じまたは紛らわしい表示を使うことから名称やブランドを守る。一方で、商標が一般名詞化すると、保護を失うことがある。

\subsubsection*{1.7.3 特許}

\addcontentsline{toc}{subsubsection}{1.7.3 特許}

特許とは、発明者が一定期間、発明を製造・販売する権利を保護する仕組みである。特許出願では、発明に至る過程を記録し、権利範囲を主張する文書を作成する。ソフトウェアに関して何が特許対象になるかは国や地域で異なる。契約中に生まれた発明の権利は、雇用主、顧客、共同所有に帰属する場合がある。

\subsubsection*{1.7.4 著作権}

\addcontentsline{toc}{subsubsection}{1.7.4 著作権}

著作権は、創作された表現を保護する仕組みである。重要なのは、著作権が「考え方そのもの」ではなく「考え方の表現」を保護する点である。ソースコード、文書、図、画面デザインの一部などは著作権と関係する可能性がある。

\subsubsection*{1.7.5 営業秘密}

\addcontentsline{toc}{subsubsection}{1.7.5 営業秘密}

営業秘密とは、一般に知られておらず、事業上の経済的価値を持ち、秘密として管理される情報である。数式、アルゴリズム、設計、手順、方法、データの集まりなどが該当し得る。営業秘密は期間制限なく保護され得るが、別の人が合法的に同じ情報を独自に発見した場合は、その人も利用できる。

\begin{center}
\small
\begin{longtable}{p{0.19\linewidth}p{0.36\linewidth}p{0.40\linewidth}}
\toprule
\textbf{知的財産の種類} & \textbf{保護する対象} & \textbf{初学者向けの見分け方} \\
\midrule
\endfirsthead
\toprule
\textbf{知的財産の種類} & \textbf{保護する対象} & \textbf{初学者向けの見分け方} \\
\midrule
\endhead
商標 & 名称、ロゴ、記号、包装など & 「誰の商品・サービスか」を示す印である。 \\
\midrule
特許 & 発明や技術的アイデア & 新しい技術的仕組みを一定期間保護する。 \\
\midrule
著作権 & 表現された作品 & ソースコードや文書など、表現された形を保護する。 \\
\midrule
営業秘密 & 秘密管理された有用情報 & 外に出さないことで価値を保つ情報である。 \\
\midrule
\bottomrule
\end{longtable}
\normalsize
\end{center}

\begin{center}
\fbox{\parbox{0.94\linewidth}{\textbf{TODO（図） 知的財産の違いを示す図を入れる。}\par\textbf{画像生成用プロンプト}\par 白背景、モノクロ、A4本文幅の4象限比較図を作成する。四つの箱として「商標」「特許」「著作権」「営業秘密」を配置し、それぞれに「保護するもの」「例」「注意点」を短く書く。中央に「知的財産」を置く。日本語ラベル、講義ノート風。}}
\end{center}

\subsubsection*{1.7.6 専門職責任}

\addcontentsline{toc}{subsubsection}{1.7.6 専門職責任}

専門職責任とは、技術者が専門家として注意義務を果たさなかった場合に問われ得る責任である。ソフトウェア技術者は、顧客や雇用主にサービスを提供するため、専門過誤、過失、能力不足の主張に備えなければならない。

責任に関する主張への防御として、標準と一般に受け入れられた実務に従っていたことを示すことが重要になる。ソフトウェアの安全性や適合性は測定が難しいため、適切な注意を払わなかったと見なされると過失の根拠になり得る。

\subsubsection*{1.7.7 法的要求}

\addcontentsline{toc}{subsubsection}{1.7.7 法的要求}

ソフトウェア技術者は、地域、国、国際的な法的枠組みの中で活動する。登録、免許、試験、教育、経験、研修などの要件、契約上の合意、責任に関する非契約上の法律などを理解する必要がある。

\subsubsection*{1.7.8 貿易遵守}

\addcontentsline{toc}{subsubsection}{1.7.8 貿易遵守}

貿易遵守とは、物品、サービス、技術の輸入、輸出、再輸出に関する法的制約を守ることである。ソフトウェア、暗号技術、クラウドサービス、開発ツール、技術文書なども対象になり得る。輸出管理、制裁対象国や企業、政府許可、輸入制限、関税などは専門性が高いため、実務では貿易遵守の専門家に確認する必要がある。

\subsubsection*{1.7.9 サイバー犯罪}

\addcontentsline{toc}{subsubsection}{1.7.9 サイバー犯罪}

サイバー犯罪とは、コンピュータ、ソフトウェア、ネットワーク、組込みソフトウェアが関わる犯罪である。コンピュータやソフトウェアが犯罪の手段になる場合も、攻撃対象になる場合もある。詐欺、不正アクセス、盗聴、スパム、脅迫、嫌がらせ、個人データや営業秘密の窃取、他のシステムへの侵入や破壊などが含まれる。

ソフトウェア技術者には、サイバー犯罪の脅威を考慮し、ソフトウェアや利用者情報を偶発的または悪意あるアクセス、利用、変更、破壊、開示から守る責任がある。

\begin{center}
\fbox{\parbox{0.94\linewidth}{\textbf{ダークパターン}\par ダークパターンとは、利用者を誤導し、利用者が望まない行動を取らせるように設計されたUI/UX上の仕掛けである。利用者の利益よりも操作・搾取を優先するため、倫理的な実務とはいえない。}}
\end{center}

\subsubsection*{1.7.10 データプライバシー}

\addcontentsline{toc}{subsubsection}{1.7.10 データプライバシー}

データプライバシーは多くの国で重要な法的要求である。SWEBOK第14章では、EUおよびEEAにおける一般データ保護規則、ならびにカリフォルニア州消費者プライバシー法が例として挙げられている。これらは、個人データの収集、利用、移転、保護、本人の権利に関する代表的な制度である。

ソフトウェア技術者は、個人データを扱う設計、ログ、分析、外部連携、クラウド利用、削除、同意、アクセス権限を、法律と倫理の両面から考える必要がある。

\subsection*{1.8 文書化}

\addcontentsline{toc}{subsection}{1.8 文書化}

文書化は、明確で、十分で、正確な情報を利害関係者へ提供する責任である。文書の適切さは、読む人のニーズによって判断される。開発者、テスト担当、運用担当、顧客、利用者、管理者は、それぞれ必要な情報が違う。

ソフトウェア技術者は、関連事実、重要なリスクとトレードオフ、ソフトウェアの使用・誤用に伴う危険な結果、ライセンスや調達元に関する情報を記録する必要がある。不適切な製品を承認すること、機密情報を開示すること、事実やデータを偽ることは避けなければならない。

\begin{center}
\small
\begin{longtable}{p{0.21\linewidth}p{0.74\linewidth}}
\toprule
\textbf{文書の読者} & \textbf{必要な情報の例} \\
\midrule
\endfirsthead
\toprule
\textbf{文書の読者} & \textbf{必要な情報の例} \\
\midrule
\endhead
開発組織内 & 要求仕様、設計文書、利用した開発ツール、テスト仕様と結果、採用した手法、開発中に生じた問題。 \\
\midrule
顧客・利用者 & そのソフトウェアがニーズを満たしそうか判断する情報、安全な使い方と危険な使い方、機密情報の保護方法、警告と重要手順。 \\
\midrule
保守担当 & 機能仕様、設計、テストスイート、必要な実行環境、変更履歴、既知の制約。 \\
\midrule
\bottomrule
\end{longtable}
\normalsize
\end{center}

\begin{center}
\fbox{\parbox{0.94\linewidth}{\textbf{保管期間}\par 文書は、少なくともソフトウェア製品のライフサイクルの間、または組織・規制上必要とされる期間、保持する必要がある。}}
\end{center}

\begin{center}
\fbox{\parbox{0.94\linewidth}{\textbf{TODO（図） 文書化の読者別情報を示す図を入れる。}\par\textbf{画像生成用プロンプト}\par 白背景、モノクロ、A4本文幅の三列図を作成する。列を「開発組織内」「顧客・利用者」「保守担当」とし、各列に必要な文書情報を短い箇条書きで配置する。中央上部に「文書化は読者ごとに必要情報が異なる」と書く。日本語ラベル、講義ノート風。}}
\end{center}

\subsection*{1.9 トレードオフ分析}

\addcontentsline{toc}{subsection}{1.9 トレードオフ分析}

トレードオフ分析とは、複数の解決案について、リスク、費用、便益を利害関係者と協力して評価し、目的に最も合う案を選ぶことである。ソフトウェアでは、性能、信頼性、安全性、セキュリティ、費用、納期、保守性、利用者体験などが競合しやすい。

プロジェクトが遅延したり予算超過したりした場合、どの要求を緩和または削除できるかを判断するためにトレードオフ分析が行われることがある。最初の手順は、設計目標を定め、その相対的な重要度を決めることである。目標の表現が曖昧だと、分析結果も曖昧になる。

トレードオフ分析は倫理的に行う必要がある。代替案を比較する基準と重みづけは、客観的かつ公平に選ぶべきである。利害の衝突がある場合は、あらかじめ開示しなければならない。

\begin{center}
\fbox{\parbox{0.94\linewidth}{\textbf{例}\par 「開発費を下げる」「信頼性を上げる」「応答時間を短くする」は同時に最大化できないことがある。トレードオフ分析では、どの目的をどれだけ重視するかを明示し、根拠を残す。}}
\end{center}

\begin{center}
\fbox{\parbox{0.94\linewidth}{\textbf{TODO（図） トレードオフ分析の流れを示す図を入れる。}\par\textbf{画像生成用プロンプト}\par 白背景、モノクロ、A4本文幅のフローチャートを作成する。「設計目標を定める」「評価基準を決める」「重みを付ける」「代替案を列挙する」「リスク・費用・便益を評価する」「利害関係者と合意する」「決定と根拠を記録する」の順に箱を並べる。重要箇所として「利益相反の開示」を点線で添える。日本語ラベル、講義ノート風。}}
\end{center}

\section*{2 集団力学と心理学}

\addcontentsline{toc}{section}{2 集団力学と心理学}

ソフトウェア工学の仕事は、多くの場合チームで行われる。チームには、開発者、テスト担当、運用担当、顧客、利用者、管理者、規制当局、外部委託先などが関わる。したがって、技術者には、他者と協力し、対立を建設的に扱い、情報を共有し、継続的に改善する力が求められる。

\subsection*{2.1 チーム・集団で働く力学}

\addcontentsline{toc}{subsection}{2.1 チーム・集団で働く力学}

成果を出すチームは、凝集性を持ち、協力的で、正直で、焦点の合った雰囲気を持つ。個人の目標とチームの目標がそろっていると、メンバーは共有成果に自然に責任と所有感を持つ。

よいチームでは、知的に正直であること、知らないことを認めること、誤りを認めること、責任・報酬・作業負荷を公平に分けることが重視される。文書、ソースコード、会話で明確に伝えることで、情報が全員に届く。レビューでは人ではなく成果物に焦点を当て、非個人的かつ建設的に行う。

ソフトウェアは多様な領域に入り込むため、技術者は多分野の環境で働けなければならない。スマートフォン、自動車、医療機器、業務システムなど、適用領域が変われば、関係者の価値観やリスクも変わる。

\subsection*{2.2 個人の認知}

\addcontentsline{toc}{subsection}{2.2 個人の認知}

個人の認知とは、人が情報を理解し、記憶し、問題を分解し、解決策を考える仕組みである。ソフトウェアは抽象度が高いため、個人の認知限界は問題解決に大きく影響する。

\begin{center}
\small
\begin{longtable}{p{0.27\linewidth}p{0.68\linewidth}}
\toprule
\textbf{認知を妨げる要因} & \textbf{説明} \\
\midrule
\endfirsthead
\toprule
\textbf{認知を妨げる要因} & \textbf{説明} \\
\midrule
\endhead
知識不足 & 対象領域、技術、標準、制約を知らないために正しく判断できない。 \\
\midrule
無意識の前提 & 自分では当然と思っている条件が、実際には成り立たない。 \\
\midrule
データ量の多さ & 情報が多すぎて、重要な点を見落とす。 \\
\midrule
失敗への恐れ & 評価や責任を恐れ、相談や試行を避ける。 \\
\midrule
文化 & 組織文化や適用領域の文化が、質問や発言を妨げる。 \\
\midrule
表現力不足 & 問題を言語化・図式化できず、共有できない。 \\
\midrule
雰囲気 & 発言しにくい職場環境が問題発見を遅らせる。 \\
\midrule
感情状態 & 疲労、不安、怒りなどが判断に影響する。 \\
\midrule
\bottomrule
\end{longtable}
\normalsize
\end{center}

\begin{center}
\fbox{\parbox{0.94\linewidth}{\textbf{知的謙虚さ}\par 知的謙虚さとは、自分の知識や判断に限界があることを認め、他者に相談し、証拠を確認する態度である。複雑なソフトウェア問題では、この態度が誤った前提を減らす。}}
\end{center}

\begin{center}
\fbox{\parbox{0.94\linewidth}{\textbf{TODO（図） 個人の認知限界と対策を示す図を入れる。}\par\textbf{画像生成用プロンプト}\par 白背景、モノクロ、A4本文幅の因果図を作成する。左に「妨げる要因」として「知識不足」「無意識の前提」「データ量」「失敗への恐れ」「文化」「感情状態」を配置し、中央に「認知負荷」を置く。右に対策として「問題分解」「集中」「相談」「読書」「試作」「レビュー」を配置する。日本語ラベル、講義ノート風。}}
\end{center}

\subsection*{2.3 問題の複雑さへの対応}

\addcontentsline{toc}{subsection}{2.3 問題の複雑さへの対応}

多くのソフトウェア工学問題は、全体を一度に扱うには複雑すぎる。そのため、チーム作業と問題分解が必要になる。問題分解とは、大きな問題を小さな部分に分け、それぞれを理解・解決し、最後に統合する考え方である。

チームでは、メンバーごとの知識、経験、視点、創造性が相補的に働く。ペアプログラミングやコードレビューは、複雑な問題に対するチームでの対処例である。

\subsection*{2.4 利害関係者との相互作用}

\addcontentsline{toc}{subsection}{2.4 利害関係者との相互作用}

ソフトウェア工学の成功は、利害関係者との良好な相互作用に依存する。利害関係者は、ソフトウェアライフサイクルのすべての段階で、支援、情報、フィードバックを提供する。初期段階では、すべての利害関係者を特定し、製品がどのように影響するかを理解することが重要である。

アジャイル開発では、利害関係者の関与が特に重要である。開発中、利害関係者は仕様や早期版へのフィードバック、優先度変更、詳細要求の明確化を提供する。保守中も、進化する要求や問題報告を通じて、ソフトウェアの拡張と改善に関わる。

\begin{center}
\fbox{\parbox{0.94\linewidth}{\textbf{TODO（図） 利害関係者関与のライフサイクル図を入れる。}\par\textbf{画像生成用プロンプト}\par 白背景、モノクロ、A4本文幅の循環図を作成する。「要求」「設計」「実装」「テスト」「運用」「保守」を円環状に配置し、外側に「利害関係者の支援・情報・フィードバック」を置く。各段階に短い例を添える。日本語ラベル、講義ノート風。}}
\end{center}

\subsection*{2.5 不確実性と曖昧さへの対応}

\addcontentsline{toc}{subsection}{2.5 不確実性と曖昧さへの対応}

ソフトウェア技術者は、不確実性と曖昧さを扱わなければならない。不確実性とは、必要な情報が不足している状態である。曖昧さとは、同じ言葉や状況が複数の意味に解釈できる状態である。

不確実性が知識不足に由来する場合は、教科書、専門誌、標準、論文、公式資料を読む、利害関係者に面談する、チームメイトや専門家に相談することで解決できることがある。容易に解消できない不確実性は、プロジェクトリスクとして扱い、見積りや価格、計画に反映する。

\begin{center}
\fbox{\parbox{0.94\linewidth}{\textbf{TODO（図） 不確実性への対応フローを入れる。}\par\textbf{画像生成用プロンプト}\par 白背景、モノクロ、A4本文幅のフローチャートを作成する。「不確実または曖昧な事項」から始め、「知識不足か？」で分岐する。Yesは「資料調査」「利害関係者へ確認」「専門家へ相談」へ進む。Noまたは未解決の場合は「リスクとして登録」「見積り・計画へ反映」「監視」へ進む。日本語ラベル、講義ノート風。}}
\end{center}

\subsection*{2.6 公平性・多様性・包摂性への対応}

\addcontentsline{toc}{subsection}{2.6 公平性・多様性・包摂性への対応}

公平性、多様性、包摂性は、チームの力学に影響する。地理的に分散したチームやコミュニケーション頻度が低いチームでは、一回一回の接触の重要性が高まる。時差があると口頭でのやり取りが減り、文化的な誤解も起こりやすい。

ソフトウェア開発では国際的な外部委託や分散開発が多いため、多文化環境は特に一般的である。プロジェクトを成功させるには、文化的・社会的規範の違いへの寛容さ、リーダーシップによる支援、より頻繁なコミュニケーション、可能な範囲での対面または同期的な会議が有効である。

性別による偏りもソフトウェア産業の課題である。より広い採用戦略、具体的で測定可能な評価基準、報酬決定手続きの透明化は、不平等を減らすための手段になり得る。

\begin{center}
\fbox{\parbox{0.94\linewidth}{\textbf{包摂性の実務的意味}\par 包摂性は「やさしい雰囲気を作る」だけではない。必要な人が議論に参加でき、情報へアクセスでき、自分の懸念を言語化でき、評価や報酬で不利にならないようにする実務上の設計である。}}
\end{center}

\section*{3 コミュニケーション技能}

\addcontentsline{toc}{section}{3 コミュニケーション技能}

ソフトウェア技術者には、口頭、読解、文章のすべてでよく伝える力が必要である。要求と納期を満たすためには、顧客、上司、同僚、供給者との明確なコミュニケーションが欠かせない。

問題解決は、情報を調べ、理解し、要約できるときに最適化される。顧客による製品受け入れと安全な利用は、適切な教育、説明、文書に依存する。技術者のキャリアも、口頭・文章での効果的な伝達能力に影響される。

\subsection*{3.1 読む・理解する・要約する}

\addcontentsline{toc}{subsection}{3.1 読む・理解する・要約する}

ソフトウェア技術者は、技術書、マニュアル、研究論文、オンライン資料、ソースコードを読んで理解できなければならない。読むことは技能向上の手段であり、工学的目標を達成するための情報収集手段でもある。

顧客が複数の解決策の中から選ぶ必要がある場合、技術者は調査結果を要約し、相手が理解できる形に整理することが求められる。ソースコードを読む力も重要である。ソフトウェアを修正、拡張、再作成するとき、技術者は提示されたコードから実装を理解し、時には設計を推測しなければならない。

\subsection*{3.2 書く}

\addcontentsline{toc}{subsection}{3.2 書く}

ソフトウェア技術者が作る文書には、ソースコード、プロジェクト計画、要求文書、リスク分析、設計文書、テスト計画、利用者マニュアル、技術報告、評価、根拠説明、図表などがある。

明確で簡潔な文章は重要である。文章は、関係者間の主要なコミュニケーション手段になることが多い。書かれた成果物は、対象読者にとってアクセス可能で、理解可能で、関連性がなければならない。

\subsection*{3.3 チームと集団のコミュニケーション}

\addcontentsline{toc}{subsection}{3.3 チームと集団のコミュニケーション}

協働するソフトウェア工学活動では、チームや集団の効果的なコミュニケーションが不可欠である。利害関係者への相談、意思決定、計画作成、問題共有は、すべてコミュニケーションを通じて行われる。

ただし、チームメンバーが増えると、可能なコミュニケーション経路の数は二次関数的に増える。たとえば、n人のチームでは、二者間の経路は n(n-1)/2 になる。人数が増えるほど、誰に何を伝えるかを設計しなければならない。

文書、電子メール、共同作業ツール、情報リポジトリは有用である。しかし、メッセージが多すぎると重要情報が埋もれる。チームによっては、対面または仮想的な同居、共用作業空間、共同レビューを重視することで、協働を促進する。

\begin{center}
\small
\begin{longtable}{p{0.33\linewidth}p{0.62\linewidth}}
\toprule
\textbf{チーム人数} & \textbf{二者間コミュニケーション経路数} \\
\midrule
\endfirsthead
\toprule
\textbf{チーム人数} & \textbf{二者間コミュニケーション経路数} \\
\midrule
\endhead
3人 & 3経路 \\
\midrule
5人 & 10経路 \\
\midrule
8人 & 28経路 \\
\midrule
10人 & 45経路 \\
\midrule
20人 & 190経路 \\
\midrule
\bottomrule
\end{longtable}
\normalsize
\end{center}

\begin{center}
\fbox{\parbox{0.94\linewidth}{\textbf{TODO（図） チーム人数とコミュニケーション経路の増加図を入れる。}\par\textbf{画像生成用プロンプト}\par 白背景、モノクロ、A4本文幅の折れ線または概念図を作成する。横軸を「チーム人数」、縦軸を「二者間コミュニケーション経路数」とし、n(n-1)/2で急増する曲線を描く。点として3人=3、5人=10、10人=45、20人=190を示す。日本語ラベル、講義ノート風。}}
\end{center}

\subsection*{3.4 プレゼンテーション技能}

\addcontentsline{toc}{subsection}{3.4 プレゼンテーション技能}

ソフトウェア技術者は、ライフサイクルの多くの場面でプレゼンテーション技能を使う。要求レビューで顧客に要求を説明する、設計レビューで構造を説明する、製品ウォークスルーを行う、利用者や運用担当へ教育する、といった場面である。

よい発表は、技術情報を一方的に話すだけではない。相手の前提知識に合わせて概念を整理し、意思決定に必要な情報を示し、質問やフィードバックを引き出す。発表で使った知識は、スライド、技術メモ、白書、ナレッジ記事などとして保存し、組織の知識資産にすることが望ましい。

\begin{center}
\fbox{\parbox{0.94\linewidth}{\textbf{TODO（図） コミュニケーション技能の関係図を入れる。}\par\textbf{画像生成用プロンプト}\par 白背景、モノクロ、A4本文幅の三角形図を作成する。三角形の頂点に「読む・理解する・要約する」「書く」「話す・発表する」を置き、中央に「相手が意思決定できるようにする」を置く。周囲に「顧客」「チーム」「運用担当」「利用者」を配置する。日本語ラベル、講義ノート風。}}
\end{center}

\section*{4 トピックと参考文献の対応}

\addcontentsline{toc}{section}{4 トピックと参考文献の対応}

この表は、SWEBOK第14章が示す「トピック対参考資料の行列」を学習用に簡略化したものである。詳しく学ぶときは、各トピックに対応する文献を読むとよい。

\begin{center}
\small
\begin{longtable}{p{0.40\linewidth}p{0.55\linewidth}}
\toprule
\textbf{トピック} & \textbf{主な参考資料} \\
\midrule
\endfirsthead
\toprule
\textbf{トピック} & \textbf{主な参考資料} \\
\midrule
\endhead
1. 専門性 & Bott et al.; Sommerville; McConnell; Tockey \\
\midrule
1.1 認定・資格証明・資格・免許 & Bott et al.; Sommerville; Washington Accord; ISO/IEC 24773-1; ISO/IEC 24773-4 \\
\midrule
1.2 倫理規程と専門職行動規範 & Bott et al.; Voland; Sommerville; McConnell; ACM Code; IEEE Code; IFIP Code; Tockey \\
\midrule
1.3 専門職団体 & Bott et al.; Sommerville; McConnell \\
\midrule
1.4 ソフトウェア工学標準 & Bott et al.; Appendix B; Sommerville; McConnell \\
\midrule
1.5 経済的影響 & Voland; Sommerville; Tockey \\
\midrule
1.6 雇用契約 & Bott et al.; ACM Code; IEEE Code; IFIP Code \\
\midrule
1.7 法的論点 & Bott et al.; Appendix B; Voland; Sommerville \\
\midrule
1.8 文書化 & Bott et al.; Voland; Sommerville; McConnell \\
\midrule
1.9 トレードオフ分析 & Voland; Sommerville; Tockey \\
\midrule
2. 集団力学と心理学 & Voland; Sommerville; McConnell; Fairley \\
\midrule
2.1 チーム・集団で働く力学 & Voland; Fairley \\
\midrule
2.2 個人の認知 & Voland; McConnell \\
\midrule
2.3 問題の複雑さへの対応 & Voland; Sommerville; McConnell \\
\midrule
2.4 利害関係者との相互作用 & Sommerville \\
\midrule
2.5 不確実性と曖昧さ & Sommerville; Fairley \\
\midrule
2.6 公平性・多様性・包摂性 & Sommerville; Fairley \\
\midrule
3. コミュニケーション技能 & Voland; Sommerville; McConnell; Fairley \\
\midrule
3.1 読む・理解する・要約する & Sommerville; McConnell \\
\midrule
3.2 書く & Voland; Sommerville \\
\midrule
3.3 チームと集団のコミュニケーション & Voland; Sommerville; McConnell; Fairley \\
\midrule
3.4 プレゼンテーション技能 & Voland; Sommerville; Fairley \\
\midrule
\bottomrule
\end{longtable}
\normalsize
\end{center}

\section*{5 発展的な読み物}

\addcontentsline{toc}{section}{5 発展的な読み物}

\begin{center}
\small
\begin{longtable}{p{0.38\linewidth}p{0.57\linewidth}}
\toprule
\textbf{文献} & \textbf{読むと得られること} \\
\midrule
\endfirsthead
\toprule
\textbf{文献} & \textbf{読むと得られること} \\
\midrule
\endhead
Weinberg, The Psychology of Computer Programming & プログラミングを個人とチームの活動として捉える古典的文献である。技術だけでなく、人の認知、チーム、組織がプログラミングに与える影響を学べる。 \\
\midrule
Kinney and Lange, Intellectual Property Law for Business Lawyers & 米国を中心に、知的財産法がどのような考え方で構成されているかを学べる。ソフトウェア契約、著作権、特許、営業秘密を考える入口になる。 \\
\midrule
\bottomrule
\end{longtable}
\normalsize
\end{center}

\section*{6 参考文献}

\addcontentsline{toc}{section}{6 参考文献}

以下はSWEBOK第14章の参考文献を、学習用に日本語資料内で参照しやすい形に整理したものである。

\begin{itemize}[leftmargin=2em]

\item F. Bott et al., Professional Issues in Software Engineering, 3rd ed., Taylor \& Francis, 2000.

\item Appendix B of the SWEBOK Guide v4.0a.

\item G. Voland, Engineering by Design, 2nd ed., Prentice-Hall, 2003.

\item I. Sommerville, Software Engineering, 10th ed., Addison-Wesley, 2016.

\item S. McConnell, Code Complete, 2nd ed., Microsoft Press, 2004.

\item 25 Years Washington Accord, IEC, 2014.

\item EUR-ACE, 2017.

\item ISO/IEC 24773-1, Software and Systems Engineering — Certification of Software and System Engineering Professionals — Part 1: General Requirements.

\item ISO/IEC 24773-4, Software and Systems Engineering — Certification of Software and Systems Engineering Professionals — Part 4: Software engineering.

\item ACM Code of Ethics and Professional Conduct, 2018.

\item IEEE Code of Ethics, 2020.

\item IFIP Code of Ethics and Professional Conduct, 2021.

\item S. Tockey, Return on Software: Maximizing the Return on Your Software Investment, Addison-Wesley, 2004.

\item R. E. Fairley, Managing and Leading Software Projects, Wiley-IEEE Computer Society Press, 2009.

\item G. M. Weinberg, The Psychology of Computer Programming: Silver Anniversary Edition, Dorset House, 1998.

\item Kinney and Lange, P.A., Intellectual Property Law for Business Lawyers, Thomson West, 2013.

\end{itemize}

\section*{7 画像 TODO 一覧}

\addcontentsline{toc}{section}{7 画像 TODO 一覧}

本文に配置した画像TODOを再掲する。実際に図を作る場合は、各TODOのプロンプトをそのまま画像生成に使うと、A4資料に合う図を作りやすい。

\begin{center}
\small
\begin{longtable}{p{0.11\linewidth}p{0.26\linewidth}p{0.62\linewidth}}
\toprule
\textbf{番号} & \textbf{TODO} & \textbf{画像生成用プロンプト} \\
\midrule
\endfirsthead
\toprule
\textbf{番号} & \textbf{TODO} & \textbf{画像生成用プロンプト} \\
\midrule
\endhead
1 & 14章の全体地図を入れる。 & 白背景、モノクロ、A4本文幅に収まる横長の学習用図解を作成する。中央に「ソフトウェア工学専門職実務」を置き、三つの大分類として「専門性」「集団力学と心理学」「コミュニケーション技能」を配置する。「専門性」から「倫理」「標準」「契約」「法的論点」「文書化」「トレードオフ分析」へ枝分かれさせる。日本語ラベル、細線、講義ノート風、余白を広めにする。 \\
\midrule
2 & 認定・資格証明・免許の違いを示す比較図を入れる。 & 白背景、モノクロ、A4本文幅の比較図を作成する。横軸に「組織を対象」から「個人を対象」、縦軸に「任意性が高い」から「法的強制力が高い」を置く。点として「認定」「資格証明」「資格」「免許」を配置し、それぞれに短い説明を添える。日本語ラベル、講義ノート風。 \\
\midrule
3 & 倫理規程違反の種類を示す図を入れる。 & 白背景、モノクロ、A4本文幅の二分割図を作成する。左に「作為による違反」、右に「不作為による違反」を配置する。左には「機密漏えい」「事実の偽造」「能力の虚偽表示」「不十分な仕事の隠蔽」を、右には「リスク未開示」「重要情報の未提供」「参考文献や貢献の未表示」「顧客利益の不十分な代表」を入れる。日本語ラベル、講義ノート風。 \\
\midrule
4 & 知的財産の違いを示す図を入れる。 & 白背景、モノクロ、A4本文幅の4象限比較図を作成する。四つの箱として「商標」「特許」「著作権」「営業秘密」を配置し、それぞれに「保護するもの」「例」「注意点」を短く書く。中央に「知的財産」を置く。日本語ラベル、講義ノート風。 \\
\midrule
5 & 文書化の読者別情報を示す図を入れる。 & 白背景、モノクロ、A4本文幅の三列図を作成する。列を「開発組織内」「顧客・利用者」「保守担当」とし、各列に必要な文書情報を短い箇条書きで配置する。中央上部に「文書化は読者ごとに必要情報が異なる」と書く。日本語ラベル、講義ノート風。 \\
\midrule
6 & トレードオフ分析の流れを示す図を入れる。 & 白背景、モノクロ、A4本文幅のフローチャートを作成する。「設計目標を定める」「評価基準を決める」「重みを付ける」「代替案を列挙する」「リスク・費用・便益を評価する」「利害関係者と合意する」「決定と根拠を記録する」の順に箱を並べる。重要箇所として「利益相反の開示」を点線で添える。日本語ラベル、講義ノート風。 \\
\midrule
7 & 個人の認知限界と対策を示す図を入れる。 & 白背景、モノクロ、A4本文幅の因果図を作成する。左に「妨げる要因」として「知識不足」「無意識の前提」「データ量」「失敗への恐れ」「文化」「感情状態」を配置し、中央に「認知負荷」を置く。右に対策として「問題分解」「集中」「相談」「読書」「試作」「レビュー」を配置する。日本語ラベル、講義ノート風。 \\
\midrule
8 & 利害関係者関与のライフサイクル図を入れる。 & 白背景、モノクロ、A4本文幅の循環図を作成する。「要求」「設計」「実装」「テスト」「運用」「保守」を円環状に配置し、外側に「利害関係者の支援・情報・フィードバック」を置く。各段階に短い例を添える。日本語ラベル、講義ノート風。 \\
\midrule
9 & 不確実性への対応フローを入れる。 & 白背景、モノクロ、A4本文幅のフローチャートを作成する。「不確実または曖昧な事項」から始め、「知識不足か？」で分岐する。Yesは「資料調査」「利害関係者へ確認」「専門家へ相談」へ進む。Noまたは未解決の場合は「リスクとして登録」「見積り・計画へ反映」「監視」へ進む。日本語ラベル、講義ノート風。 \\
\midrule
10 & チーム人数とコミュニケーション経路の増加図を入れる。 & 白背景、モノクロ、A4本文幅の折れ線または概念図を作成する。横軸を「チーム人数」、縦軸を「二者間コミュニケーション経路数」とし、n(n-1)/2で急増する曲線を描く。点として3人=3、5人=10、10人=45、20人=190を示す。日本語ラベル、講義ノート風。 \\
\midrule
11 & コミュニケーション技能の関係図を入れる。 & 白背景、モノクロ、A4本文幅の三角形図を作成する。三角形の頂点に「読む・理解する・要約する」「書く」「話す・発表する」を置き、中央に「相手が意思決定できるようにする」を置く。周囲に「顧客」「チーム」「運用担当」「利用者」を配置する。日本語ラベル、講義ノート風。 \\
\midrule
\bottomrule
\end{longtable}
\normalsize
\end{center}

\section*{8 章末確認問題}

\addcontentsline{toc}{section}{8 章末確認問題}

\begin{itemize}[leftmargin=2em]

\item 専門職実務とは何か。単なるプログラミング能力とどう違うかを説明せよ。

\item 認定、資格証明、資格、免許の違いを、対象と強制力に注目して説明せよ。

\item 倫理規程違反が作為だけでなく不作為でも起こる理由を説明せよ。

\item 専門職団体がソフトウェア技術者にとって重要な理由を三つ挙げよ。

\item 標準に従うことが、品質面と法的面で役立つ理由を説明せよ。

\item 秘密保持契約と知的財産契約で確認すべき点を述べよ。

\item 商標、特許、著作権、営業秘密の違いを一つずつ例を挙げて説明せよ。

\item トレードオフ分析で、評価基準と重みづけを明示する必要がある理由を説明せよ。

\item チーム人数が増えるとコミュニケーションが難しくなる理由を、経路数の観点から説明せよ。

\item 不確実性と曖昧さの違いを説明し、それぞれへの対処を述べよ。

\item 公平性・多様性・包摂性が、分散ソフトウェア開発で重要になる理由を説明せよ。

\item 「読む・理解する・要約する」「書く」「発表する」が、ソフトウェア技術者の仕事にどう関係するかを説明せよ。

\end{itemize}

\begin{center}
\fbox{\parbox{0.94\linewidth}{\textbf{解答の方向性}\par 専門職実務は、知識、技能、態度、倫理、標準、説明責任を含む。認定は教育機関やプログラム、資格証明と資格は個人の能力、免許は法的な業務権限に注目する。標準は最低限の共通土台と説明根拠になる。トレードオフ分析では、目的、評価基準、重み、利害関係者との合意、利益相反の開示が重要である。}}
\end{center}

\end{document}